\section{近期研究（2026 年 8--11 月）}
\subsection{街景前馈重建与几何一致性}

环视前馈重建的难点分布在三个位置：输入特征要提供几何先验，预测时要把原语分配到有用的表面，训练时还要把相机外观和几何分开。以下方法分别从这三个位置切入。

\paragraph{VGGD：把基础模型的几何先验放到环视前馈网络前端}
2026.08~\cite{2608_10682}。单时刻环视图像的相机间重叠很少，几何不确定性很难只靠增加解码器容量解决。VGGD 从预训练 VGGT 提取多视图几何特征，再分别处理几何一致性与外观补全，并在训练早期渐进引入尺度约束。预测头结合像素和体积两种参数化，输出用于新视图渲染的高斯。

论文在 nuScenes 单帧协议中评价画质，并以相对于 Depth Anything V2 伪深度的 Pearson 相关评价结构。该相关指标对全局平移和正比例缩放不敏感，因此支持相对深度一致性的比较，无法单独证明米制深度准确。

\paragraph{LGS：用删点和加点干预学习高斯密度控制}
2026.08~\cite{2608_11077}。LiDAR 引导的前馈重建常让每个观测点初始化一个高斯，重建质量却未必随初始化密度均匀分布。LGS 在训练阶段模拟增加或删除基元，用局部的梯度响应作为结构调整的训练信号，使推断时能够按需要改变高斯数量。

论文在 Waymo 与 PandaSet 上评价，并将动态区域和整张图的画质及深度误差分别列出。与逐场景增密相比，该方法将增密决策转化为训练得到的前馈策略，其在不同扫描模式或数据域上的表现需要重新验证。

\paragraph{SG-GS：用语义区域和边界调节城市高斯分配}
2026.10~\cite{2610_16292}。复杂街景里统一的增密阈值可能忽略细小构件和边界。SG-GS 将 SAM 分割得到的语义区域引入高斯密度控制，在不同区域动态分配梯度阈值，并在结构边界加强增密。该方法也用语义约束组织优化损失，使结构细节与原语数量的分配相互关联。

在 Derrick 数据上，相同划分和训练设置下，作者比较原始 3DGS 并报告 PSNR 提升 1.34 dB。语义分区质量决定阈值能否覆盖真实细部，渲染指标还需与几何误差或下游驾驶指标分别验证。

\paragraph{MCSplat：把环视外观差异从几何监督中分离}
2026.10~\cite{2610_25345}。车载相机的曝光和图像处理差异会让同一表面在不同相机里呈现不同颜色。若重建模型直接降低观测图像之间的光度误差，几何可能被迫吸收这部分变化。MCSplat 利用深度引导的跨相机对应构建统一外观，并用多尺度双边网格 Transformer 把统一外观的渲染结果恢复到相机各自的观测外观。这个设计保留从真实图像得到的训练信号，在外观映射中解释相机差异。

Waymo 协议以时刻 t 的图像评价输入视图重建，以 t+1 的图像评价新视角合成。跨视图深度和法向一致性另行评价。作者也在 nuScenes 测试迁移表现，并指出训练集外观有限时，少数场景中恢复的相机域外观仍会偏色。跨车道外推示例提供定性证据。

\subsection{相机观测退化与天气}

恢复受天气或挡风玻璃影响的观测，需要先区分污染附着在哪个坐标系。可以绑定在相机，也可以随场景深度变化；仿真则要控制这些效应何时进入渲染。

\paragraph{GSRAIN：在街景高斯中分别校准雨滴与整体雨效}
2026.08~\cite{2608_02177}。街景渲染的雨效同时包含快速变化的局部雨滴和低频的整体外观。GSRAIN 根据实测降雨数据构建高频雨滴模型，用受几何引导的单步扩散模型生成低频雨效，再把两者融合进同一高斯场景。作者给出的雨强控制区间为 0 至 13 毫米每小时。

论文的总体雨效对比在 Waymo 上用 8.5 毫米每小时的模拟设定，报告 FID 149.09，对照 CycleGAN-Turbo 与 WeatherEdit 分别为 155.71 和 157.94。整体雨强可控范围与这组图像质量比较的具体雨强属于不同条件。检测和闭环结果反映所测场景、雨量与驾驶模型的组合。

\paragraph{DerainSplat：从稀疏雨天观测恢复可渲染的干净场景}
2026.08~\cite{2608_02191}。与在已有高斯场景中施加降雨的任务相比，DerainSplat 从少量受天气污染的图像出发，把天气因素与场景分开估计。天气网络预测阴天光照、随深度变化的雾、雨线和镜头雨滴，由此生成两类支持图。一类调整跨视图代价体的匹配，另一类在深度对齐后融合用于修复外观的证据。低几何支持的位置还会削弱高斯不透明度，以抑制伪结构。

论文对合成雨天输入做定量比较，也在真实雨天 HydroViews 场景上展示无微调结果。真实雨天没有与之对齐的干净目标，因此这部分主要用定性结果评价；它与合成雨天的成对指标应分开解读。

\paragraph{RIDE-Splat：在高斯重建之外解释挡风玻璃反射}
2026.10~\cite{2610_34632}。透过挡风玻璃拍摄时，外部街景随相机运动产生视差，座舱反射却大体跟着相机。RIDE-Splat 先联合估计世界坐标下的三维高斯和一张相机坐标下共享的有符号 RGB 图，让后者解释反射增亮与透射衰减。随后去掉反射成分，再用处理后的观测拟合干净场景。

推断时只渲染去反射的高斯场景。评价区分纯合成、真实场景加合成反射和真实挡风玻璃反射。真实的无反射参考来自另一次采集，时序和环境难以完全对齐；强烈独立运动的反射也超出共享图近似的适用范围。

\subsection{轨迹外推与视频修复}

高斯渲染在日志轨迹附近通常有较多测量支撑，轨迹外推后可见范围改变，修复模块需限制对可靠原图内容的改写。

\paragraph{SPVC：按空间条件修复整段驾驶渲染}
2026.08~\cite{2608_17420}。新轨迹和车辆插入会同时破坏背景结构、前景外观和帧间连续性。SPVC 在视频级处理这些伪影，并用相机位姿、三维框和高精地图约束修复内容。训练样本用欠约束渲染及插入车辆产生的退化构造，让模型同时见到背景扭曲与前景接缝。

作者在 nuScenes、Waymo 和 PandaSet 上检查新视角渲染修复及资产插入伪影。位姿、三维框或地图不准会影响结构约束；修复后图像的观感还需与原始几何和下游任务的一致性分开评价。

\paragraph{PRG-Fusion：按重建证据决定保留、修复或生成}
2026.09~\cite{2609_06948}。沿日志以外的轨迹渲染驾驶视频时，同一图像可能同时包含可信的重建表面、受退化影响但仍有测量支撑的区域，以及完全缺少观测的区域。PRG-Fusion 根据三维高斯渲染和几何证据按像素划分这三类区域。可信部分保留高斯外观，中间部分用 LiDAR 投影约束结构，缺少支撑的部分交给视频先验补全。

作者在 Waymo 验证集选取八个含较多动态交通参与者的片段进行评价。方法依赖预先重建与对齐的 LiDAR；比较轨迹外推时，应分别报告可由重建或测量支持的区域，以及由生成模型补齐的区域。

\subsection{动态实例与闭环仿真}

可独立编辑的交通参与者需要同时具备可见表面、实例布局和可控运动。实例重建、不可见面补全和行为合成承担不同的工作。

\paragraph{DecoupleGS：把动态交通参与者做成可重用资产}
2026.08~\cite{2608_01761}。闭环测试需要按规划器输出调整自车和周边车辆的轨迹。DecoupleGS 将静态街景与规范坐标系下的动态交通参与者分别表示。动态资产通过感知裁剪和向量量化控制渲染成本，跟随地图拓扑进行几何注册，再以代理重光照减小资产和背景的外观差异。

实验以 nuScenes 和 PandaSet 片段作为背景，从 3DRealCar 选取可操纵的车辆资产，分别评估渲染效率、几何与光度一致性、日志回放及闭环测试。代理重光照是效率优先的近似，地图注册依赖语义拓扑，当前实验主要覆盖刚性交通参与者和相机仿真。

\paragraph{USR-Drive：共同去噪动态高斯与实例三维框}
2026.08~\cite{2608_19036}。驾驶重建输出稠密、可渲染的高斯，感知输出可以描述实例的三维框。USR-Drive 在有位姿环视视频的条件下联合生成两类三维状态。高斯向框提供稠密几何线索，框向动态高斯提供实例级布局约束。

实验采用 nuScenes 六路环视相机的八帧片段，并在 VKitti 的 400 个样本上做零样本评价。位姿用于统一位置编码的米制锚定，不能把这一条件描述为任意无标定视频。渲染重建和三维检测仍按不同指标评价。

\paragraph{单侧车辆观测下的米制初始化与对称补全}
2026.09~\cite{2609_08841}。单侧稀疏图像用于车辆资产时，米制尺度和未观测侧外观是两个不同的问题。该方法用视觉基础模型预测米制点云作为高斯中心锚点，通过学习得到的编码和解码模块估计高斯属性，再在高斯空间沿车辆的近似对称面复制可靠原语。

论文在 3DRealCar 用同侧四张图构建资产，并在车辆对侧八个视角评价。与 RECAST 的学习式未见面生成相比，这条路线依靠车辆的双侧对称约束，车门、标识和附件等非对称细节仍需观测。

\paragraph{RECAST：给车辆资产补齐日志中没有看到的视角}
2026.09~\cite{2609_31374}。即使把动态车辆从街景中独立出来，日志重建的车辆仍可能缺少背面或侧面。RECAST 从一张已分割的车辆图像生成可从各方向渲染的高斯，再将资产注册到日志重建的场景。作者构建 RECAR 车辆数据适配单图到三维生成先验，并通过修改车辆和自车轨迹验证闭环渲染。

闭环比较从九个 Waymo 场景选取一个目标车辆并移除其他动态参与者，使用固定的 GTRS-Dense 规划器。结果侧重目标车辆超出日志轨迹时的渲染和规划反应，复杂多参与者交互还需独立评估。

\paragraph{ControlPed：把风险轨迹转化为多视角可渲染行人}
2026.10~\cite{2610_06171}。端到端驾驶评估需要同时控制行人进入道路的位置、时刻和外观。ControlPed 将交通冲突轨迹转成受文本与轨迹约束的人体运动，再驱动高斯化身渲染多相机观测。高斯化身承担三维外观和动画，风险条件来自交通轨迹与姿态生成流程。

评估保留 88 个经检查并渲染的场景，覆盖七个端到端模型；报告的平均 HDScore 从 88.8 降到 47.4。当前设置为开环，因此这一差值应解读为在固定渲染观测下的退化，不能作为交互闭环成绩。
